\section{Availability and Limitations}

The artifact repository contains the \tool Python package, adapter scripts,
metric runners, generated CSV, Markdown, and \LaTeX{} tables, provenance logs,
the BibTeX/source audit used for this draft, an MIT license, and a reviewer
quickstart. The reviewer entry point is the tag
\texttt{audit-sid-cikm-resource-v0.1}, which exposes a clean-checkout verifier
for the published tables and figure. Large local artifacts are kept out of git
and described through manifests so that reviewers can distinguish public
reproducibility assets from local caches.

\begin{table}[t]
\caption{Reviewer-facing artifact checklist. Columns indicate whether each
asset is inspectable in a clean checkout, executable without private caches,
and tied to a PDF claim.}
\label{tab:artifact-package}
\small
\setlength{\tabcolsep}{3pt}
\begin{tabular}{@{}lccc@{}}
\toprule
Asset group & Inspect & Run & PDF claim \\
\midrule
Quickstart/license & yes & -- & package \\
Unit tests/verifier & yes & yes & code path \\
Figure/table CSVs & yes & yes & Tables 1--3 \\
Manifest/claim audit & yes & -- & row roles \\
Extra CSV panels & yes & -- & support \\
\bottomrule
\end{tabular}
\end{table}

The evidence supports a conservative resource-demo claim. The toolkit does not
reproduce TIGER end-to-end; the GRID Musical row is a controlled feature-text
export. The ReSID Sports balanced GAOQ path was stopped after constrained
k-means became CPU-bound, so the main ReSID evidence is the bounded Musical
export. CARD remains a fallback/stressor path, not faithful CARD
reproduction~\cite{wei2026card}. D2 includes a bounded interaction-qualified
controller but not causal harm; D3 is not proven monotonic with Recall@K or
NDCG; D5a is structural rather than measured latency; and D6 is optional drift
evidence.

These boundaries define the resource's scope. \tool audits tokenizer
artifacts; it is not a new \sid tokenizer, a complete
generative-recommender evaluation, or an industrial online-impact study.
Artifact diagnostics can reveal pressure points before generator training, but
they do not replace ranking evaluation. Future versions should add causal
collision-harm validation, generator-output diagnostics, stronger same-dataset
method coverage, and unified search/recommendation checks~\cite{penha2025jointsid}.
Industrial SID deployment studies further motivate treating these artifact
properties as engineering objects~\cite{ju2026snapchatsid}.

\paragraph{Reviewer workflow and claim discipline.}
Reviewers inspect the manifest to distinguish named-method, controlled,
optional, and literature-only rows, run the clean-checkout verifier, and compare
published CSVs with the paper claims. Full metric rebuilds are optional and
require local caches. The resource separates \emph{mapping statements} directly
supported by SID-assignment tables, \emph{diagnostic statements} that join
mappings with metadata or interactions, and \emph{system-quality statements}
that require generator training, ranking evaluation, or serving traces. This
paper makes the first two kinds of statements and explicitly not the third.

\paragraph{Maintenance plan.}
New adapters emit the same normalized tables, declare their evidence role, and
record feature or checkpoint provenance. Versioned manifests and fixed-name
``latest'' reports preserve stable entry points and timestamped provenance. The
Musical case study is therefore not a ranking of GRID against ReSID; it shows
how one artifact contract turns collisions, collaborative prefix recovery, tail
allocation, and prefix fan-out into separately checkable review objects.
